\documentclass[11pt,a4paper]{article}

\usepackage[utf8]{inputenc}
\usepackage[T1]{fontenc}
\usepackage[french]{babel}
\usepackage{lmodern}
\usepackage{geometry}
\geometry{margin=2.2cm}
\usepackage{amsmath,amssymb}
\usepackage{booktabs}
\usepackage{graphicx}
\usepackage{hyperref}
\usepackage{listings}
\usepackage{xcolor}
\usepackage{caption}

\hypersetup{
  colorlinks=true,
  linkcolor=blue!50!black,
  citecolor=blue!50!black,
  urlcolor=blue!50!black
}

\lstset{
  basicstyle=\ttfamily\footnotesize,
  breaklines=true,
  frame=single,
  columns=fullflexible,
  showstringspaces=false
}

\title{PredLDA\\[0.4em]
\large Package R d'analyse discriminante linéaire prédictive\\[0.3em]
\normalsize M2 SISE --- Programmation R --- 2026--2027}
\author{
Thomas Fonferrier \and Ryan Ducret \and Aya Saadi \and Xilian Jing\\
Équipe G
}
\date{\today}

\begin{document}
\maketitle

\begin{abstract}
Ce rapport présente \textbf{PredLDA}, un package R proposant l'analyse
discriminante linéaire (LDA) prédictive via une classe R6 d'interface proche de
scikit-learn (\texttt{fit}, \texttt{predict}, \texttt{predict\_proba}). Nous
décrivons les formules, l'architecture logicielle, le pseudo-code des
algorithmes, l'application Shiny de démonstration, les tests de validation et
les limites actuelles du cœur numérique.
\end{abstract}

\tableofcontents
\vspace{1em}

%----------------------------------------------------------------------
\section{Introduction et objectifs}
%----------------------------------------------------------------------

L'objectif du projet est de livrer un package R installable depuis GitHub,
accompagné d'une aide aux normes R (anglais), d'un tutoriel reproductible, d'une
application Shiny et du présent rapport. Le cœur statistique implémente la LDA
prédictive pour une cible catégorielle à \(K \ge 2\) classes, en exposant les
fonctions de classification linéaires.

Le dépôt GitHub est
\url{https://github.com/thomas-fonferrier/projet-lda-r}
(sous-dossier installable \texttt{PredLDA/}).

%----------------------------------------------------------------------
\section{Rappels méthodologiques : LDA prédictive}
%----------------------------------------------------------------------

\subsection{Hypothèses et notation}

Soient \(X \in \mathbb{R}^{n \times p}\) les variables explicatives numériques
et \(y\) une variable cible à \(K\) modalités. Pour chaque classe \(k\), on note
\(\mu_k\) la moyenne empirique, \(\pi_k\) la probabilité a priori, et \(W\) la
matrice de covariance intra-classe \emph{poolée}. Sous l'hypothèse de
covariances égales entre classes, la règle de Bayes conduit à des scores
linéaires \cite{hastie2009esl}.

\subsection{Fonctions de classification}

Pour chaque classe \(k\) :
\begin{align}
  a_k &= W^{-1} \mu_k, \\
  a_{0k} &= \log \pi_k - \tfrac{1}{2} \mu_k^\top W^{-1} \mu_k, \\
  d_k(x) &= a_{0k} + a_k^\top x.
\end{align}
La classe prédite est \(\arg\max_k d_k(x)\). Les probabilités a posteriori sont
obtenues par softmax des scores \(d_k(x)\).

\subsection{Covariance poolée}

Après centrage de chaque observation sur la moyenne de sa classe
(\(X_c\)), on estime
\begin{equation}
  W = \frac{X_c^\top X_c}{n - K}.
\end{equation}
Si \(W\) est singulière (variables colinéaires ou constantes), une
pseudo-inverse de Moore--Penrose est utilisée et un avertissement est émis.

\subsection{Priors}

Par défaut, \(\pi_k = n_k / n\). L'utilisateur peut fournir un vecteur de
priors strictement positifs (renormalisé). Les priors n'affectent que les
intercepts via \(\log \pi_k\).

\subsection{Cas binaire}

Pour \(K=2\), la comparaison de \(d_1\) et \(d_2\) est équivalente à
\(D(x)=d_1(x)-d_2(x)\). La méthode reste identique au cas multi-classes ;
\(D(x)\) n'est qu'une réécriture pédagogique affichée par \texttt{summary()}.

%----------------------------------------------------------------------
\section{Architecture du package}
%----------------------------------------------------------------------

\subsection{Organisation des modules}

\begin{table}[ht]
\centering
\begin{tabular}{@{}ll@{}}
\toprule
Élément & Rôle \\
\midrule
\texttt{R/stats\_utils.R} & Briques statistiques internes \\
\texttt{R/lda\_r6.R} & Classe publique \texttt{LDA} (R6) \\
\texttt{man/} & Aide Rd générée par roxygen2 \\
\texttt{tests/} & Tests unitaires et référence \texttt{MASS::lda} \\
\texttt{inst/extdata/} & Jeux d'exemple du tutoriel \\
\texttt{inst/shiny/} & Application Shiny embarquée \\
\texttt{vignettes/} & Vignette \texttt{predlda-intro} \\
\texttt{report/} & Sources LaTeX de ce rapport \\
\bottomrule
\end{tabular}
\caption{Architecture du package PredLDA.}
\end{table}

La séparation volontaire est : \emph{calculs purs} dans
\texttt{stats\_utils.R}, \emph{validation / API / stockage} dans la classe R6.
L'application Shiny n'implémente pas l'algorithme ; elle appelle les méthodes
publiques du package.

\subsection{Interface publique (R6)}

Alignement sur scikit-learn :
\begin{itemize}
  \item \texttt{LDA\$new(priors, tol)} --- constructeur ;
  \item \texttt{\$fit(X, y)} --- apprentissage ;
  \item \texttt{\$predict(X)} / \texttt{\$predict\_proba(X)} --- déploiement ;
  \item \texttt{\$score(X, y)} --- exactitude ;
  \item \texttt{\$print()} / \texttt{\$summary()} --- restitution.
\end{itemize}
Propriétés exposées après \texttt{fit} : \texttt{coefficients},
\texttt{means}, \texttt{priors}, \texttt{within\_cov}, \texttt{classes},
\texttt{feature\_names}, etc.

\subsection{Contrats d'entrée}

\begin{itemize}
  \item \(X\) : data.frame ou matrice, colonnes numériques finies, sans NA ;
  \item \(y\) : au moins deux classes ; \(n > K\) ;
  \item à la prédiction, alignement des colonnes \textbf{par nom}.
\end{itemize}

%----------------------------------------------------------------------
\section{Pseudo-code des algorithmes}
%----------------------------------------------------------------------

\subsection{Apprentissage (\texttt{fit})}

\begin{lstlisting}
FIT(X, y; priors_user, tol):
  X <- preparer_et_valider(X)          # numerique, noms uniques
  y <- factoriser_et_valider(y)
  counts <- effectifs(y)
  M <- moyennes_par_classe(X, y, counts)
  W <- covariance_poolée(X, y, M)      # Xc'Xc / (n-K)
  (W_inv, pseudo) <- inverser(W, tol)  # Cholesky ou Moore-Penrose
  pi <- resoudre_priors(priors_user, counts)
  pour chaque classe k:
    a_k  <- W_inv * mu_k
    a0_k <- log(pi_k) - 0.5 * mu_k' * W_inv * mu_k
  stocker coefficients, means, priors, W
  retourner self
\end{lstlisting}

\subsection{Prédiction}

\begin{lstlisting}
DECISION_FUNCTION(X):
  X <- aligner_colonnes_par_nom(X, feature_names)
  S <- X %*% A + a0                    # scores d_k(x)
  retourner S

PREDICT(X):
  retourner classes[argmax_ligne(DECISION_FUNCTION(X))]

PREDICT_PROBA(X):
  retourner softmax_lignes(DECISION_FUNCTION(X))
\end{lstlisting}

\subsection{Softmax numériquement stable}

Pour chaque ligne de scores \(s\), on calcule
\(\exp(s_j - \max s) / \sum_j \exp(s_j - \max s)\) afin d'éviter les
débordements numériques.

%----------------------------------------------------------------------
\section{Choix d'implémentation}
%----------------------------------------------------------------------

\begin{itemize}
  \item \textbf{Inversion de \(W\)} : tentative Cholesky (\texttt{chol} /
    \texttt{chol2inv}) ; repli eigen/pseudo-inverse si matrice mal
    conditionnée.
  \item \textbf{API chaînable} : \texttt{fit} retourne \texttt{self} de façon
    invisible (\texttt{LDA\$new()\$fit(X,y)\$predict(X)}).
  \item \textbf{Robustesse prédiction} : matching par nom, colonnes
    supplémentaires ignorées, vecteur nommé accepté pour une observation.
  \item \textbf{Pas de dépendance algorithmique lourde} : seul \texttt{R6} est
    en \texttt{Imports} ; \texttt{MASS} sert aux tests de référence.
\end{itemize}

%----------------------------------------------------------------------
\section{Application Shiny}
%----------------------------------------------------------------------

L'application (\texttt{inst/shiny/app.R}) permet de :
\begin{enumerate}
  \item importer des données (iris enrichi, csv/tsv/txt/xlsx) ;
  \item choisir cible, variables explicatives numériques et variables
        illustratives (exclues du modèle) ;
  \item entraîner via \texttt{LDA\$fit} ;
  \item évaluer (sur apprentissage ou split stratifié 80/20, graine 2026) ;
  \item prédire un individu ou un fichier, avec exports CSV.
\end{enumerate}
Les métriques affichées (exactitude, précision, rappel, F1, F1 macro,
exactitude équilibrée) sont descriptives ; elles n'appartiennent pas au cœur
statistique minimal mais documentent l'usage du modèle.

%----------------------------------------------------------------------
\section{Tests et validation}
%----------------------------------------------------------------------

Deux batteries principales :
\begin{itemize}
  \item \texttt{tests/test-LDA.R} : API, validations, cas singulier, accord
        avec \texttt{MASS::lda} sur iris ;
  \item \texttt{tests/test-reference.R} : coefficients confrontés à une
        formule \og textbook \fg{} indépendante ; jeux iris, mtcars,
        ToothGrowth, Pima, fgl, crabs ; prédiction hors échantillon Pima.
\end{itemize}
L'accord des postérieurs avec \texttt{MASS::lda} \cite{venables2002mass} à une
tolérance de \(10^{-8}\) constitue la validation externe principale.

%----------------------------------------------------------------------
\section{Installation et documentation}
%----------------------------------------------------------------------

\begin{lstlisting}
remotes::install_github(
  "thomas-fonferrier/projet-lda-r",
  subdir = "PredLDA",
  upgrade = "never"
)
\end{lstlisting}

Documentation utilisateur :
\begin{itemize}
  \item pages d'aide \texttt{?LDA} (anglais, normes R / roxygen2) ;
  \item README et vignette \texttt{predlda-intro} (tutoriel reproductible) ;
  \item fichiers d'exemple dans \texttt{inst/extdata/}.
\end{itemize}

%----------------------------------------------------------------------
\section{Limites et perspectives}
%----------------------------------------------------------------------

Le cœur actuel exige des prédicteurs numériques. Parmi les extensions prévues
par le cahier des charges et non finalisées dans le moteur : prétraitement des
explicatives catégorielles, stratégie automatique de sélection de variables,
indicateurs de pertinence des variables. L'interface Shiny rend ces limites
explicites et n'assimile pas la sélection manuelle à une procédure statistique
automatique.

%----------------------------------------------------------------------
\section{Conclusion}
%----------------------------------------------------------------------

PredLDA fournit une LDA prédictive installable, documentée et testée, avec une
API R6 claire et une application Shiny de démonstration. L'implémentation suit
les formules classiques \cite{hastie2009esl} et se valide contre
\texttt{MASS::lda}.

%----------------------------------------------------------------------
\bibliographystyle{plain}
\bibliography{references}

\appendix
\section{Répartition indicative des modules}
Les calculs de \texttt{stats\_utils.R} sont internes (non exportés). Seule la
classe \texttt{LDA} est exportée dans le \texttt{NAMESPACE}.

\section{Jeux d'exemple livrés}
\texttt{iris\_apprentissage.tsv}, \texttt{iris\_apprentissage\_fr.csv},
\texttt{iris\_nouveaux.tsv} (également utilisés par l'application Shiny).

\end{document}
